\chapter{Tổng kết}
\label{ch:tongket}

\section{Tổng hợp phát hiện}

Cả ba câu hỏi kiểm định đều chỉ về một hướng: con số mAP 0,97 cần được nhìn thận trọng hơn.
Nó có thể được data leakage nâng đỡ (CH-A), đồng thời che lấp chênh lệch lớn giữa các lớp
(CH-B), và chênh lệch đó bắt nguồn từ loại vật thể chứ không phải mất cân bằng (CH-C).

\begin{itemize}
  \item \textbf{CH-A --- Rò rỉ:} 202 ảnh trùng chính xác và 423 ảnh gần trùng xuyên split chưa
  được khử ở tập đánh giá.
  \item \textbf{CH-B --- Khoảng cách:} đèn 0,541 so với biển 0,854 (chênh 0,313; $p = 0{,}0095$).
  \item \textbf{CH-C --- Nguyên nhân:} độ hiếm không liên quan ($\rho = 0{,}032$); lỗi tập trung
  ở đèn tín hiệu.
\end{itemize}

\section{Khuyến nghị}

\begin{itemize}
  \item Khử trùng lặp xuyên split ở \emph{cả} tập valid/test trước khi báo cáo mAP; báo cáo thêm
  mAP sau khi loại near-duplicate để thấy mức chênh.
  \item Báo cáo AP theo từng lớp và theo nhóm (đèn/biển) thay vì chỉ mAP tổng, để không che giấu
  điểm yếu.
  \item Cải thiện lớp đèn tín hiệu: tăng độ phân giải đầu vào (\texttt{imgsz} 896/1024), bổ sung
  mẫu đèn khó, và phân tích ma trận nhầm lẫn đỏ$\leftrightarrow$xanh.
  \item Không dựa vào oversampling lớp hiếm để sửa lỗi đèn, vì độ hiếm không phải nguyên nhân.
\end{itemize}

\section{Hạn chế}

Đề tài ở chế độ \textit{results-only} nên không chạy lại mô hình sau khi khử rò rỉ để đo trực
tiếp mức sụt mAP; kết luận về ảnh hưởng của rò rỉ mang tính định lượng phơi nhiễm, không phải đo
can thiệp. Nhóm đèn chỉ có hai lớp nên kiểm định thống kê tuy có ý nghĩa vẫn cần thận trọng khi
khái quát.

\section{Sản phẩm kèm theo}

\begin{itemize}
  \item Pipeline phân tích tái lập (\texttt{src/}) kèm kiểm thử tự động (\texttt{pytest}, 6/6 đạt).
  \item Ứng dụng web Streamlit (\texttt{app/app.py}) trình bày toàn bộ phân tích.
  \item Báo cáo Word (\texttt{.docx}), báo cáo LaTeX (mẫu trường) và slide trình bày.
\end{itemize}
